\documentclass[10.5pt]{article}
\usepackage[a4paper,top=0.6in,bottom=0.9in,left=0.9in,right=0.9in]{geometry}
\usepackage{titlesec}
\usepackage{enumitem}
\usepackage[hidelinks]{hyperref}
\usepackage{parskip}
\usepackage{titling}

%----- Section formatting -----
\titleformat{\section}{\bfseries\large}{}{0em}{}[\titlerule]
\titleformat{\subsection}{\bfseries\normalsize}{}{0em}{}
\titlespacing*{\section}{0pt}{0.9em}{0.5em}

% bullet spacing
\setlist[itemize]{itemsep=1.5pt, topsep=1pt, partopsep=0pt, parsep=0pt, leftmargin=1.4em}

%----- Header -----
\renewcommand{\maketitle}{
  \begin{center}
    {\LARGE \textbf{Radha Pingle}}\\[0.3em]
    \rule{0.5\textwidth}{0.4pt}\\[0.4em]
    \href{mailto:rpingle@andrew.cmu.edu}{rpingle@andrew.cmu.edu} \,$\bullet$\, 973-981-7814
  \end{center}
  \vspace{-1.0em}
}

\begin{document}
\maketitle

%====================================================
\vspace{-0.5em}
\section{Education}
\textbf{Carnegie Mellon University}, Pittsburgh, PA \hfill \textit{May 2028} \\
Bachelor of Science, Cognitive Science + Statistics and Data
Science\\ 
\textbf{Nominated \& Selected as Quantitative Social Science Scholar (QSSS Program)} \\[0.7em]
\textbf{Relevant Coursework:} Methods for Statistics \& Data Science, Applied Quantitative Social Science Seminar I \& II, Intro to Research Methods, Principles of Computing, Fundamentals of Programming \& Computer Science, Principles of Imperative Computation, Cognitive Psychology


%====================================================
\section{Experience}

\noindent\textbf{Data Science / Product Intern – Epistemix} \hfill \textit{May 2025 – Aug 2025}\\[-1.3em]
\begin{itemize}
  \item Owned the 0-to-1 development of ``Public Health Copilot,'' an LLM tool (Gemini 1.5 Flash) delivering ZIP-code-specific insights, iterating on response quality through structured prompt changes and testing.
  \item Drove product decisions by deploying 15 surveys and conducting 7 user interviews, translating findings into concrete feature and prompt-design changes.
  \item Built a mock ZIP-code dataset to test model behavior and validate response consistency before release.
  \item Developed a React/JavaScript conversational interface and integrated it with the model via API calls.
  \item Translated user research findings into product improvements across prompt design and the front-end experience.
\end{itemize}

\noindent\textbf{Research Assistant – Cognitive \& Social Development Lab, CMU} \hfill \textit{May 2025 – Aug 2025}\\[-1.3em]
\begin{itemize}
  \item Reviewed dozens of academic papers on how visual illustration styles shape children's perception and implicit bias, surfacing key research gaps to inform product direction.
  \item Designed an original experiment and hypotheses testing how facial depiction style influences race-based perception, scoping the problem from an open research question.
\end{itemize}

\noindent\textbf{Website Manager – Professor Emerita, CMU} \hfill \textit{Nov 2024 – Aug 2025}\\[-1.3em]
\begin{itemize}
  \item Partnered with a faculty stakeholder to redesign and update her academic website, prioritizing changes that improved content organization, navigation, and overall clarity.
  \item Managed site structure and clean version updates using FileZilla, VSCode, and SFTP.
\end{itemize}

\noindent\textbf{Machine Learning RA – Laboratory for Computational Perception} \hfill 
\textit{May 2026 – Present}\\[-1.3em]
\begin{itemize}
  \item Scaled a labeled natural-sound dataset by over 260\% (275 to 1,000+ validated samples) to power human and AI auditory-categorization research.
  \item Used Python to curate, validate, and organize short audio clips with category labels.
  \item Support dataset quality control by reviewing sound clips, refining labels, and preparing data for human--model comparison experiments.
\end{itemize}

\noindent\textbf{Cognitive Development Researcher -- Kid Neuro Lab} \hfill
\textit{May 2026 -- Present}\\[-1.3em]
\begin{itemize}
  \item Processed structural and functional fMRI data for 5 participants across 3 studies investigating STEM skill development in children.
  \item Conducted quality control checks on fMRI images, identifying motion artifacts and other scan-quality issues prior to analysis.
  \item Assisted with fMRI data collection for child participants, supporting scan setup, participant flow, and experimental procedures.
\end{itemize}


%====================================================
\section{School Projects}
\begin{itemize}
  \item \textbf{Principles of Computing (15-110)} – Trained a language model on Grimm \& Andersen fairy tales and evaluated stylistic patterns through statistical frequency analysis. 
  \item \textbf{Fundamentals of Programming and Computer Science (15-112)} – Built a Python game with randomized recipes, priority-based scheduling, and sequencing-dependent scoring logic.
\end{itemize}
%====================================================
\section{Technical Skills}
Python, R, SQL, React, JavaScript, HTML, Java, C++, C, Git, Figma, Matlab

\end{document}
